\documentclass{ximera}

\begin{document}

    \title{In-Class Assignment 5}
    \begin{abstract} Solve the following problems. \end{abstract}
    \maketitle

\begin{exercise}
Given a series $\sum\limits_{n=1}^\infty a_n$ with $a_n \neq 0$, the Ratio Test states that if $\{a_n\}$ satisfies
\begin{align*}
    \lim \left|\frac{a_{n+1}}{a_n}\right| = r < 1,
\end{align*}
then the series converges absolutely.
\begin{enumerate}
    \item Let $r'$ satisfy $r< r'<1$. (Why must such an $r'$ exist?)  Explain why there exists an $N$ such that $n \geq N$ implies $|a_{n+1}| \leq |a_n|r'$.
    \item Why does $|a_N|\sum (r')^n$ necessarily converge?
    \item Now, show that $\sum |a_n|$ converges.
\end{enumerate}
\end{exercise}

\newpage

\begin{exercise}
    Show that if  $a_n > 0$ and $\lim (na_n) = l$, $l \neq 0$, then the series $\sum a_n$ diverges.
\end{exercise}
\vfill

\begin{exercise}
    Assume $a_n > 0$ and $\lim (n^2a_n) $ exists.  Show that $\sum a_n$ converges.
\end{exercise}
\vfill

\end{document}